\subsubsection{The Dark Dimension and Majorana Neutrinos --- arXiv:2609.10668}

\textbf{Authors:} Arturo de Giorgi, Dhruv Pasari

\paragraph{主な主張}
$0.1$--$10\,\mu\mathrm{m}$規模のDark Dimensionにbulk Majoranaニュートリノを導入した模型について、Daya Bay・KATRIN・KamLAND-Zenの制限を比較し、想定する質量領域ではKATRINが最も強くYukawa結合を$\mathcal{O}(10^{-3})$まで制限することを示した\cite{deGiorgi:2026DarkDimensionMajorana}。

\paragraph{新規性}
全KK towerを含むニュートリノを放出しない二重ベータ崩壊の有効Majorana質量を閉じた形で求め、核内運動量がコンパクト化スケールとMajorana質量を上回る場合にKK状態の寄与が打ち消し合って崩壊振幅が抑制されることを示した。

\paragraph{位置付け}
bulk Majorana質量を持つ余剰次元ニュートリノ模型を複数の地上実験で検証する研究であり、$T^2$など異なるコンパクト化幾何におけるKKスペクトルと実験制約を比較する際の基準となる。
